% from ega4-10-fr.tex
\oldpage[IV-3]{100}
Pour les deux premières, ce n'est autre que la définition (10.2.2)~; comme la
topologie de $X$ est l'image réciproque par $f$ de celle de $f(X)$ les cinq
dernières applications, où l'on remplace $Y$ par $f(X)$, sont bijectives. On
peut donc (par (10.2.1, \emph{b)})) se borner au cas où $X$ est un sous-espace
très dense de $Y$, et le fait que les applications
$\mathfrak{Lf}(Y)\to\mathfrak{Lf}(X)$,
$\mathfrak{Qc}(Y)\to\mathfrak{Qc}(X)$,
$\mathfrak{Lqc}(Y)\to\mathfrak{Lqc}(X)$ sont bijectives a déjà été prouvé
(10.1.2). Toute partie localement constructible étant localement
quasi-constructible, les applications $\mathfrak C(Y)\to\mathfrak C(X)$ et
$\mathfrak{Lc}(Y)\to\mathfrak{Lc}(X)$ sont injectives~; en outre, pour tout
ouvert $U\subset Y$, la restriction $f^{-1}(U)\to U$ de $f$ est un
quasi-homéomorphisme (10.2.4), donc si l'on montre que
$\mathfrak C(Y)\to\mathfrak C(X)$ est surjective, il en sera de même de
$\mathfrak{Lc}(Y)\to\mathfrak{Lc}(X)$. Mais en vertu de (10.2.6), toute partie
ouverte rétrocompacte $Z$ dans $X$ est de la forme $f^{-1}(T)$, où $T$ est
ouvert rétrocompact dans $Y$~; cela prouve évidemment la surjectivité de
$\mathfrak C(Y)\to\mathfrak C(X)$.

\begin{remarks}[10.2.8]
\label{IV.10.2.8-fr}
\upshape (i) On a vu dans la démonstration de (10.2.1) que si $f$ est un
quasi-homéomorphisme, l'application canonique
\begin{equation}
\label{IV.10.2.8.1-fr}
\tag{10.2.8.1}
\Gamma(Y,\mathcal F)\to\Gamma(X,f^*(\mathcal F))
\end{equation}
est un isomorphisme de groupes abéliens fonctoriel en $\mathcal F$ dans la
catégorie des faisceaux de groupes abéliens sur $Y$. Comme $f^*$ est exact
dans cette catégorie ($\mathbf 0_{\mathrm I}$, 3.7.2), cela implique que les
homomorphismes canoniques (T, 3.2.2)
\[
\mathrm H^i(Y,\mathcal F)\xrightarrow{\sim}
\mathrm H^i(X,f^*(\mathcal F))
\]
sont \emph{bijectifs} pour tout $i$.

(ii) Si $f:X\to Y$ est une application continue et $\mathcal F$,
$\mathcal G$ deux faisceaux de groupes abéliens sur $Y$, on a
\[
f^*(\mathcal F\otimes_{\mathbf Z_Y}\mathcal G)
=f^*(\mathcal F)\otimes_{\mathbf Z_X}f^*(\mathcal G)
\]
à un isomorphisme canonique près ($\mathbf 0_{\mathrm I}$, 4.3.3). On en
conclut que si $f$ est un quasi-homéomorphisme, $f^*$ est encore une
\emph{équivalence} de la catégorie des faisceaux d'anneaux sur $Y$, et de la
catégorie des faisceaux d'anneaux sur $X$~; la donnée d'une structure d'espace
annelé sur $Y$ est donc équivalente à celle d'une structure d'espace annelé
sur $X$.

Étant donnés deux espaces annelés $(X,\mathcal O_X)$, $(Y,\mathcal O_Y)$, nous
dirons qu'un morphisme d'espaces annelés
$f=(\psi,\theta):(X,\mathcal O_X)\to(Y,\mathcal O_Y)$ est un
\emph{quasi-isomorphisme} si $\psi$ est un quasi-homéomorphisme de $X$ dans
$Y$ et $\theta^\sharp:\psi^*(\mathcal O_Y)\to\mathcal O_X$ un isomorphisme de
faisceaux d'anneaux~; lorsqu'il en est ainsi, l'espace annelé
$(X,\mathcal O_X)$ est entièrement déterminé, à isomorphisme près, par
$(Y,\mathcal O_Y)$, l'espace $X$ et le quasi-homéomorphisme $\psi$ (que l'on
peut prendre arbitraire). Lorsque $f$ est un quasi-isomorphisme d'espaces
annelés, le foncteur
\[
\mathcal F\mapsto f^*(\mathcal F)
\]
est une \emph{équivalence} de la catégorie des $\mathcal O_Y$-Modules et de
celle des $\mathcal O_X$-Modules, puisque $f^*(\mathcal F)$ s'identifie
canoniquement ici à $\psi^*(\mathcal F)$. On en conclut par exemple des
isomorphismes de bi-$\partial$-foncteurs
\[
\operatorname{Ext}^i_{\mathcal O_Y}(\mathcal F,\mathcal G)
\xrightarrow{\sim}
\operatorname{Ext}^i_{\mathcal O_X}(f^*(\mathcal F),f^*(\mathcal G)).
\]
De façon générale, on peut dire que les constructions usuelles de la théorie
des faisceaux et de l'algèbre homologique, faites sur l'espace annelé $Y$ ou
l'espace annelé $X$, sont équivalentes.
\end{remarks}

